%% Master CV - Rafael González Seliva
%% Comprehensive overview of all competencies, experience, and achievements.
%% Use as a master reference when tailoring role-specific CVs with /apply.
%%
%% Compile with: lualatex main_example.tex
%% (pdflatex often fails on modern MiKTeX with fontawesome5 font-expansion errors.)

\documentclass[11pt,a4paper,sans]{moderncv}
\moderncvstyle{banking}
\moderncvcolor{blue}

% Force the name and section headings to render in moderncv blue (color1).
\renewcommand*{\namefont}{\fontsize{34}{36}\bfseries\upshape}
\colorlet{firstnamecolor}{color1}
\colorlet{lastnamecolor}{color1}
\colorlet{namecolor}{color1}
\renewcommand*{\sectionstyle}[1]{{\sectionfont\color{color1}#1}}

\usepackage[utf8]{inputenc}
\ifpdftex\usepackage[T1]{fontenc}\fi
\AtEndPreamble{\hypersetup{
    colorlinks=true,
    linkcolor=blue,
    filecolor=magenta,
    urlcolor=blue,
    pdftitle={Rafael González Seliva - CV},
    pdfpagemode=UseNone,
}}
\usepackage[scale=0.80]{geometry}
\usepackage{import}
\usepackage{needspace}

% Personal data
\name{Rafael}{González Seliva}
\address{Terrassa, Barcelona}{}{}
\phone[mobile]{+34 677 420 408}
\email{gonzalezselivarafael@gmail.com}
\extrainfo{\href{https://www.linkedin.com/in/rafael-gonzalez-seliva}{LinkedIn}}

\begin{document}

\makecvtitle

% ============================================================
%     PROFILE STATEMENT
% ============================================================

\vspace{6pt}
\small{Senior Frontend Developer con más de 6 años de experiencia en Angular, especializado en diseñar la arquitectura y liderar el desarrollo de aplicaciones en sectores como el aeronáutico, ferroviario, movilidad eléctrica y salud. Experiencia con Angular moderno (Signals, Zoneless), CI/CD, testing automatizado con Jest y Playwright, y librerías de componentes en Storybook. Integro la IA en el flujo de desarrollo con Claude Code, agentes, MCP y gestión de contexto.}

% ============================================================
%     CORE COMPETENCIES
% ============================================================

\section{Competencias Clave}
\vspace{1pt}
\begin{itemize}

\item \textbf{Frontend Angular}: Angular (hasta v20), TypeScript, RxJS, HTML, SCSS, Signals, Zoneless, SSR, Standalone components, ChangeDetectionStrategy.OnPush, lazy loading

\item \textbf{Testing}: Playwright (e2e, Chrome DevTools Protocol), Jest, Jasmine, Karma — cobertura \>80\% en proyectos de producción

\item \textbf{Tooling \& CI/CD}: GitLab CI/CD, Jenkins, Docker, Nx, Storybook

\item \textbf{IA aplicada al desarrollo}: Claude Code, agentes, MCP, gestión eficiente de contexto para automatizar flujos de desarrollo

\item \textbf{Cloud \& Backend}: AWS (S3, CloudFront, Lambda, EC2, DynamoDB); PHP, MySQL (experiencia inicial)

\end{itemize}

% ============================================================
%     PROFESSIONAL EXPERIENCE
% ============================================================

\section{Experiencia Profesional}
\vspace{3pt}
\begin{itemize}

\item{\cventry{abr 2023-Actualidad}{Desarrollador Frontend}{Circontrol}{Viladecavalls, Barcelona}{}{\vspace{1pt}
\begin{itemize}
    \item {Diseñé la arquitectura y desarrollé varios proyectos para administrar instalaciones de cargadores de coches eléctricos con Angular 20, usando librerías gráficas como PrimeNG y Foblex.}
    \item {Reduje los tiempos de carga de una aplicación en más de un 50\% al actualizar Angular y aplicar prácticas de sus guías de estilo: lazy loading, ChangeDetectionStrategy.OnPush, Zoneless y Signals.}
    \item {Optimicé el rendimiento de un editor de diagramas interactivos (Foblex) de más de 1000 nodos: en tests de estrés automatizados con Playwright y Chrome DevTools Protocol, el tiempo de carga se redujo un 91\%, el DOM un 55\% y seleccionar un nodo pasó de 107 ms a menos de 1 ms.}
    \item {Implanté pipelines CI/CD en GitLab con Docker, logrando un proceso de actualización un 30\% más rápido.}
    \item {Integré Playwright y Jest para pruebas automatizadas, alcanzando más del 80\% de cobertura de tests unitarios.}
    \item {Creé un documento que unifica las convenciones de desarrollo del equipo de frontend para aplicarlas de forma consistente a lo largo de los proyectos.}
    \item {Implementé agentes de IA con Claude Code para automatizar flujos de desarrollo, optimizando el consumo de tokens mediante gestión eficiente del contexto.}
\end{itemize}}}

\vspace{3pt}

\needspace{5\baselineskip}
\item{\cventry{oct 2021-abr 2023}{Desarrollador Frontend}{Capgemini}{Barcelona}{}{\vspace{1pt}
\begin{itemize}
    \item {\textbf{Proyecto Airbus} (dic 2022 - abr 2023): Desarrollé una aplicación con Angular, integrando Jenkins para CI/CD y PrimeNG para la interfaz, siguiendo metodologías ágiles.}
    \item {\textbf{Proyecto ADIF} (dic 2021 - dic 2022): Lideré un equipo frontend de 4-5 personas utilizando Angular Material como librería principal de componentes, mejorando la consistencia visual y reduciendo el tiempo de desarrollo.}
    \item {\textbf{Proyecto Storybook interno} (oct - dic 2021): Desarrollé una librería de unos 50 componentes Angular en Storybook para varios proyectos internos, unificando la base de código y acelerando la entrega de productos.}
\end{itemize}}}

\vspace{3pt}

\needspace{5\baselineskip}
\item{\cventry{oct 2018-feb 2021}{Desarrollador Full Stack}{Aslogic}{Rubí, Barcelona}{}{\vspace{1pt}
\begin{itemize}
    \item {Desarrollé aplicaciones web del ámbito sanitario usando jQuery, Angular, Bootstrap, PHP y MySQL.}
\end{itemize}}}

\end{itemize}

% ============================================================
%     EDUCATION
% ============================================================

\section{Educación}
\vspace{1pt}
\begin{itemize}

\item{\cventry{2017-2020}{CFGS Desarrollo de Aplicaciones Web}{INS Nicolau Copèrnic}{Terrassa, España}{}{}}

\vspace{3pt}

\item{\cventry{2017-2019}{CFGS Desarrollo de Aplicaciones Multiplataforma}{INS Nicolau Copèrnic}{Terrassa, España}{}{}}

\end{itemize}

% ============================================================
%     LANGUAGES
% ============================================================

\section{Idiomas}
\vspace{1pt}
\begin{itemize}
\item {Español (nativo), Catalán (nativo), Inglés (B2 certificado — Professional Working).}
\end{itemize}

% ============================================================
%     REFERENCES
% ============================================================

\section{Referencias}
\vspace{1pt}
\begin{itemize}
\item{Disponibles a solicitud.}
\end{itemize}

\end{document}
